\section{Kaggle lab: distance before and after scaling}
\label{sec:kaggle-breast-geometry}

\begin{kaggleproject}{Breast Cancer Wisconsin --- geometry and similarity}
\kagglesource{https://www.kaggle.com/datasets/uciml/breast-cancer-wisconsin-data}{Breast Cancer Wisconsin (Diagnostic) Data Set}
\textbf{File.} \texttt{all\_dataset/breast\_cancer\_dataset/data.csv}\\
\textbf{Question.} How much does feature scale change Euclidean geometry, and how is cosine similarity asking a different geometric question?\\
\textbf{Before running.} Look at the feature names and predict why raw Euclidean distance may be dominated by a few measurement scales.

\begin{labmode}
This is a guided mathematical experiment. Use the reference code to verify the geometry, then change only the selected features so you can explain the effect of scale.
\end{labmode}

\lstinputlisting[style=mlpython]{all_experiments/lab04-breast_cancer_geometry.py}

\textbf{Interpretation.} The point is not that standardized Euclidean distance is always correct. It is that a distance-based algorithm silently inherits every unit and scale choice in the representation.

\begin{tryit}
Choose three features measured on noticeably different scales. Compute the pairwise distances using only those columns before and after scaling. Which feature appears to dominate the raw distance, and why?
\end{tryit}
\end{kaggleproject}

\begin{labdeliverable}
Submit one notebook with a small table of raw versus standardized Euclidean distances for three deliberately different-scale features, one cosine-similarity comparison, and a short explanation of which representation choice changed the geometry and why.
\end{labdeliverable}
